% ===
\section{Introduction}
% ===

The relationship between economic activity and the natural environment has long been a subject of scholarly and social inquiry, particularly with respect to the tension between the productive use of natural resources and ecological sustainability. At the theoretical level, this tension reflects a deeper limitation in conventional production theory: although land, capital, and labor are commonly treated as fundamental inputs of production, the ecological systems that sustain them are often modeled as external constraints, fixed stocks, or sources of externalities, rather than as endogenous and dynamic components of the production process \cite{jehle2011, stiglitz1989reflections, dasgupta2001human}. As a result, environmental degradation is often treated as an ex post consequence of production decisions, to be addressed through corrective mechanisms such as taxation, regulation, property-rights allocation, or collective governance arrangements \cite{pigou2010liberty, hardin1968tragedy, ostrom2002commonpool}.

In practical terms, this theoretical separation is reflected in legal and institutional arrangements that impose preservation duties on private agents without fully integrating ecological preservation into their economic rationale. For example, an economic agent may be required to preserve part of a biome even when doing so reduces the amount of productive land available and, consequently, constrains the agent’s production capacity, as occurs in Brazil, where legislation restricts land use to varying degrees according to the biome \cite{brasil2012florestal}. This requirement creates an opportunity cost for agricultural production, since land that could otherwise be used productively must be preserved in the interest of protecting a collective good. The agent therefore faces a trade-off between production and preservation, often framed as a conflict between economic growth and environmental conservation.

This dichotomy can be understood through three recurring paradigms of human--environment interaction: predation, mutualism, and extractivism. Predation denotes the unsustainable and often illegal depletion of natural resources; mutualism refers here to a preservation-oriented ideal in which human use is strongly limited in favor of ecological integrity; and extractivism occupies an intermediate position, allowing resource use within supposedly sustainable or regenerative limits \cite{gudynas2009diez, acosta2013extractivism}. From an efficiency perspective, however, none of these paradigms fully resolves the production--preservation trade-off. Predation disregards ecological constraints and generates negative externalities, while strict mutualism may excessively restrict productive potential \cite{dasgupta2001human, ostrom1990commons}. Extractivism, in turn, may delay ecological collapse but remains vulnerable to regulatory failure: if too restrictive, it underuses productive resources; if too permissive, it reproduces the overuse problem it was meant to avoid \cite{chichilnisky2017north}. Therefore, natural resource allocation tends to deviate from Pareto efficiency whenever private incentives are not aligned with social and ecological costs, which motivates the search for institutional arrangements capable of reconciling productive efficiency with environmental sustainability \cite{hanley2007environmental}.

Against this background, the central problem addressed in this article is whether environmental preservation can be theoretically reframed not merely as a constraint on production, but as an endogenous component of economic decision-making. More specifically, the article asks whether private productivity and profitability can be reconciled with socially desirable environmental preservation through incentive-compatible arrangements capable of overcoming the limitations of the prevailing extractivist logic.

The hypothesis advanced here is that, when ecological variables are incorporated into production theory and when incentives are designed to reflect spatial, intertemporal, and behavioral constraints, environmental preservation may become compatible with individual rationality rather than dependent solely on coercion or moral commitment. The objective of the article is therefore to develop the conceptual foundations of economic commensalism as a framework for endogenizing nature in production theory and for reinterpreting preservation as a potentially rational and productivity-enhancing economic strategy.

This article argues that such antagonism may not be inherent to economic activity itself, but rather to the way economic theory and environmental institutions have historically conceptualized nature: not as a productive factor to be strategically managed, but as an external limit to be corrected, protected, or compensated. Against this background, the article proposes a framework of economic commensalism, in which nature is endogenized within production theory and environmental preservation is examined as a potentially rational, incentive-compatible, and productivity-enhancing strategy.

























\newpage

% fim %